% mode: ini
% Aimed at section 1375, "Display math should end with $$", at its three sites: a one-$ display, the
% $ that ends a number, and the alignment's in align.rs; insert_dollar_sign, "Missing $ inserted" from \par
% and each command not allowed in the mode; and \eqno and \leqno, start_eq_no, where \eqno is refused:
% in text math, in a group, in a number already, in a \vcenter, outside math.
\input prelude
\def\e#1{\vrule width#1pt height1pt}
\def\sh{\showbox1 }
\long\def\t#1{\setbox1\hbox{#1}\sh}
\long\def\d#1{\setbox1\vbox{#1}\sh}
% One $ ends a display: the next token is put back, the paragraph resumes.
\d{\noindent$$a$ b\par}
\d{\noindent$$a$\par}
\d{\noindent$$a$\relax b\par}
\d{\noindent$$a\eqno b$ c\par}
\d{\noindent$$a\eqno b$\par}
\d{\noindent$$a\leqno b$\relax c\par}
\d{\noindent$$a\eqno b$$ c\par}
% Macro expansion finds the second $.  A \let of $ is a $ too; \relax is not.
\def\dl{$}
\let\mdl=$
\d{\noindent$$a$\dl b\par}
\d{\noindent$$a\mdl\mdl b\par}
\d{\noindent\mdl\mdl a$$ b\par}
\d{\noindent$\mdl a$$ b\par}
\d{\noindent$\relax$ a\par}
\d{\noindent$\ifx a a$\fi a$$ b\par}
% get_token is what init_math reads with, so \ifmmode is not expanded and it is not a display:
\d{\noindent$\ifmmode a\fi$\par}
% $$ in other modes: in internal vertical mode it begins a paragraph; in \hbox it is two.
\d{$$a$$ b\par}
\t{$$a$$}
\t{a$$b}
\t{a$$$$b}
% Missing $: \par in a formula, in a group in a formula, and in a number.
\d{\noindent$$a\par b$$\par}
\d{\noindent$$a\eqno b\par c$$\par}
\t{$a\par b}
\setbox1\hbox{${a\par b}\sh
\t{$\left(a\par b}
\d{\noindent$\mathchoice{a\par}{}{}{}$\par}
\t{$\vcenter{\par}$}
% vmode things in a formula, in restricted horizontal mode and in a paragraph.
\t{$\hrule}
\t{$\vskip 1pt}
\t{$\vfil}
\t{$\vss}
\t{$\vfilneg}
\t{$\unvbox0}
\t{$\unvcopy0}
\d{\noindent$\hrule\par}
\d{\noindent$\vskip 1pt\par}
\d{\noindent$\vfil\par}
\d{\noindent$\unvbox0\par}
\d{\noindent$$a\hrule\par}
\t{$\moveleft 1pt\hbox{}$}
\t{$\moveright 1pt\hbox{}$}
\t{$\halign{#\cr a\cr}$}
\t{${\halign{#\cr a\cr}}$}
\t{$\noindent a$}
\t{$\indent a$}
\t{$\unhbox0$}
\t{$\unhcopy0$}
\t{$\accent65 a$}
\t{$\ a$}
\t{$/\/a$}
\t{$\raise 1pt\hbox{}$}
\t{$\lower 1pt\hbox{}$}
\t{$\special{x}\write1{y}\openout1=z \closeout1 \immediate\closeout1$}
\t{$\vadjust{}\insert100{}\mark{m}\penalty 1 \discretionary{}{}{}$}
% math things in text, and in a vertical list: each opens a formula.
\t{^a$}
\t{_a$}
\t{\mathchar"141$}
\t{\mathord a$}
\t{\mathaccent"141 a$}
\t{\mathchoice{}{}{}{}$}
\t{\underline a$}
\t{\overline a$}
\t{\radical0 a$}
\t{\displaystyle$}
\t{\textstyle$}
\t{\scriptstyle$}
\t{\scriptscriptstyle$}
\t{\nonscript$}
\t{\mkern 1mu$}
\t{\mskip 1mu$}
\t{\vcenter{}$}
\t{\left($}
\t{\right)$}
\t{\over$}
\t{\atop$}
\t{\above 1pt$}
\t{\limits$}
\t{\nolimits$}
\t{\delimiter0 $}
\d{^a$\par}
\d{\mathchar"141$\par}
\d{\vcenter{}$\par}
\d{\eqno$\par}
% \eqno: refused unless privileged in math_shift_group.
\t{$a\eqno b$}
\t{$a\leqno b$}
\d{\noindent$$a\eqno b\eqno c$$\par}
\d{\noindent$$a\eqno b\leqno c$$\par}
\d{\noindent$${a\eqno b}$$\par}
\d{\noindent$$\left(a\eqno b\right)$$\par}
\d{\noindent$$\mathchoice{a\eqno b}{}{}{}$$\par}
\d{\noindent$$\vcenter{\hbox{$a\eqno b$}\eqno}$$\par}
\d{\noindent$$\vcenter{\eqno}$$\par}
\d{\noindent$$\begingroup a\eqno b\endgroup$$\par}
\d{\noindent$$a\over b\eqno c\over d$$\par}
% A group that \eqno leaves: \tracinggroups is e-TeX, so these show by the save stack's use, with
% \eqno's and \leqno's own \everymath; "saved(0)" is what chooses \leqno.
\everymath={\showlists}
\d{\noindent$$a\eqno b$$\par}
\d{\noindent$$a\leqno b$$\par}
\everymath={}
% \mathsurround's nodes around a number are none, and \fam comes back.
\d{\noindent$$\fam=1 a\eqno\showthe\fam\fam=2 \showthe\fam b$$\par}
\d{\noindent$$\showthe\fam a\fam=3 \eqno\showthe\fam b$$\par\showthe\fam}
% From here on, each error's recovery may leave a list open, so these run deeper in the nest.
\t{$\valign{#\cr\hbox{a}\cr}$}
\t{\eqno$}
\t{\leqno$}
% The last: \end in a formula, in an \hbox.  It ends the formula, then the box, then the job.
\t{$\end$}
